\documentclass[11pt]{article}

\usepackage[margin=1.1in]{geometry}
\usepackage{amsmath}
\usepackage{amssymb}
\usepackage{booktabs}
\usepackage{microtype}
\usepackage{xcolor}
\usepackage{caption}
\usepackage{array}
\usepackage[colorlinks=true, linkcolor=blue!50!black, citecolor=blue!50!black, urlcolor=blue!50!black]{hyperref}

\captionsetup{font=small, labelfont=bf}

\newcommand{\code}[1]{\texttt{#1}}
\newcommand{\usff}{\code{us54}}

\title{\textbf{What a Public Log Reveals:\\ Observing Play Style in Literature (Canadian Fish)}}
\author{Allen Hsieh\\[2pt] {\small FishAI project}}
\date{August 2026}

\begin{document}

\maketitle

\begin{abstract}
\noindent
In the \usff{} dialect of Literature (Canadian Fish), every action is broadcast: each ask and
its result, each declaration with the true holders of all six cards, and the public hand
counts. This paper studies the \emph{observation} problem in isolation from the adaptation
problem treated by its companion: given only that public channel, how much of a seat's play
style can be recovered from a single game, and what fundamentally limits the recovery? We
specify the channel exactly --- including what it certifies (a hit locates a card; an ask
certifies a licence; a declaration reveals its true holders, making foreign and own-hand-only
declares \emph{exact} public observations, not inferences) and what it structurally omits
(declining to declare emits no event, so declare-window patience is invisible to any log-based
observer). From the channel we build a 14-dimensional feature vector of rates and shares per
seat, calibrate a deliberately simple diagonal-Gaussian classifier on 1{,}350 mirror games,
and report two small empirical methods findings: the Gaussian normaliser term must be dropped
(it handed one style 40 of 60 control-game reads through a calibration-tightness bonus), and
per-style variances must be shrunk halfway to pooled. In the calibrated observable space the
roster's geometry is stark: four of the nine styles form a quadrangle nearly at a single point
(pairwise pooled-$z$ distance $0.20$--$0.25$ --- the companion papers' inert-axis finding
arriving in observable space), while the Turtle sits $3.2$--$3.7$ from that cluster.
Measured on 1{,}800 held-out cross-play games, single-game top-1 accuracy over nine styles
rises from $16.4\%$ at 40 events to $22.4\%$ at the full log against an $11.1\%$ chance rate
--- twice chance, and no better than that --- with per-style accuracy from $50.4\%$ (Ghost)
down to $5.6\%$ (Balanced). The ceiling is not the classifier's: styles that diverge from the
control on $0.39$--$2.89\%$ of decisions --- roughly one to three policy-revealing decisions
per seat per game, some of them invisible declines --- upper-bound what \emph{any} observer of
this channel can recover in one game. We close by locating this instrument inside the
opponent-detection program of the project's research synthesis: which panel statistics this
implementation ships, and which remain future work.
\end{abstract}

\section{Introduction}
\label{sec:intro}

Literature (Canadian Fish) is a six-player, two-team card game in which all hidden information
originates at the deal and every subsequent action is public. That makes it an unusually clean
domain for a question that is usually confounded: \emph{how much of an opponent's temperament
is visible in what they publicly do?} In most imperfect-information games the observer's
problem is entangled with hidden state --- one cannot tell a bold player from a lucky one
without knowing their cards. In Literature the log eventually tells you the cards too: every
successful ask moves a named card face up, and every declaration reveals, by rule, the true
holder of all six cards it names.

The FishAI project's v0.5 paper \cite{v05} measured \emph{which} of nine parameterised play
styles wins, holding skill fixed by giving every style one identical full-strength inference
engine. Its forthcoming v1.0 companion \cite{v10} asks whether an agent that identifies its
opponents' styles online and adapts can beat the fixed champion. Between those two questions
sits a third, which this paper isolates: the observation problem itself. Before asking whether
a read is \emph{useful}, one should measure whether a read is \emph{possible} --- and the
answer turns out to have enough structure, and enough honest negative content, to warrant its
own report.

Concretely, this paper documents the observation layer shipped in the repository as
\code{observe.ts} and \code{classify.ts} \cite{repo}: a single pass over the public log that
produces a fixed feature vector per seat, and a diagonal-Gaussian classifier over those
vectors calibrated offline from mirror games. Its contributions are:

\begin{enumerate}
\item \textbf{An exact specification of the observation channel}
  (Section~\ref{sec:channel}): what the \usff{} public log certifies --- including the
  result that foreign declares and own-hand-only declares are \emph{exact} public
  observations because declarations reveal true holders --- and what it structurally cannot
  show, chiefly that declining to declare emits no event.
\item \textbf{A 14-feature behavioural fingerprint} (Section~\ref{sec:features}), every entry
  a rate or share, with the measured game-length confound that forced that normalisation.
\item \textbf{Two empirical methods findings from calibration}
  (Section~\ref{sec:classifier}): dropping the Gaussian normaliser, and shrinking per-style
  variances halfway to pooled. Both were made by measurement, not taste, and both are the
  kind of small correction that vanishes from most write-ups.
\item \textbf{The geometry of the roster in observable space}
  (Section~\ref{sec:separability}): a pooled-$z$ distance matrix computed from the committed
  fingerprints, in which the inert-axis finding of the companion papers
  \cite{inertpaper,v05} reappears as four styles collapsed onto nearly a single point.
\item \textbf{Measured single-game accuracy and its confusion structure}
  (Section~\ref{sec:accuracy}): $22.4\%$ top-1 over nine styles against $11.1\%$ chance on
  held-out cross-play games, with the Ghost and Scout acting as attractors that receive
  $47.4\%$ of all reads.
\item \textbf{The fundamental limit and the program} (Section~\ref{sec:limit}): decision
  divergence rates of $0.39$--$2.89\%$ \cite{styles,v05} upper-bound any log-based observer,
  and the instrument is located inside the opponent-detection panel of the project's research
  synthesis \cite{playstyles}, with the shipped and unshipped panel rows stated explicitly.
\end{enumerate}

Throughout, ``the classifier'' means the committed implementation with the committed
calibration \cite{fingerprints,repo}; every quantitative claim traces to a committed artifact
\cite{artifact,fingerprints} or to the repository's measurement documents \cite{rules,styles}.

\section{The observation channel}
\label{sec:channel}

\subsection{What the public log contains}

The \usff{} engine's public log \cite{rules,repo} is a sequence of typed events:
\code{game\_started}; \code{ask} with asker, target, the named card, and whether it hit;
\code{claim} (a declaration) with the claimer, the book, the claimer's stated assignments,
the \emph{true} holders of all six cards, and the outcome; and the bookkeeping events
\code{pass}, \code{designate}, \code{player\_out}, \code{endgame}, \code{game\_over}. Hand
counts are public at all times. This is exactly the payload a human at the table would have:
bots consume only a \code{SeatView}, and a test enforces that no hidden hand is reachable
\cite{v05,repo}.

\subsection{What the log certifies exactly}
\label{sec:exact}

The observation layer records only facts the rules force to be true; nothing in the feature
extractor is an estimate. Four certifications carry the whole construction \cite{rules}:

\begin{enumerate}
\item \textbf{A hit locates a card.} The named card transfers face up
  (\code{RULES\_US54.md} row 9), so after a hit the card is publicly at the asker until it
  moves again or its book resolves. This yields a public card-location map.
\item \textbf{An ask certifies a licence.} Holding at least one card of the asked book is the
  only licence to ask into it (row 6), so every ask proves the asker held such a card at that
  moment --- and row 7 proves the asker \emph{lacked} the named card. The certification is
  decayed conservatively: it is dropped when the book resolves, when the seat runs out of
  cards, and when a hit strips a card of that book from the seat, since the stripped card may
  have been the licence itself.
\item \textbf{A declaration reveals its true holders.} The \code{claim} event carries the
  actual holder of each of the six cards, whatever the outcome. Consequently a \emph{foreign}
  declare (the claimer holds no card of the book: the claimer appears nowhere among the true
  holders) and an \emph{own-hand-only} declare (every true holder is the claimer) are
  \textbf{exact public observations}. These are precisely the public signatures of the
  \code{foreignDeclare} and \code{declareOnlyOwnHand} style mechanisms of the roster
  \cite{styles,v05}, and their exactness is what makes the corresponding features noiseless.
\item \textbf{Hand counts replay.} Counts are re-derived from the log rather than read from
  the view: every seat starts at the dealt hand size, a hit moves one card, and a claim
  removes each of the six cards from whoever the reveal says held it. This is deliberate:
  the classifier evaluates \emph{truncated} logs (Section~\ref{sec:classifier}), and a
  truncated view's own top-level counts describe the end of the game, not the truncation
  point. The replayed counts are correct at every prefix, and the test suite pins them
  against the engine's counts at full length and at mid-game snapshots \cite{repo}.
\end{enumerate}

\subsection{What the log cannot show}
\label{sec:blind}

Declining to declare emits \textbf{no event}: the \usff{} declare window advances silently
through seats that pass on their option \cite{rules,repo}. Declare-window \emph{patience} ---
how long a seat sat on a declarable book before firing --- is therefore invisible to any
observer of this channel, not merely to this implementation. No shipped feature depends on
it. The feature that comes closest, \code{declareBackload}, measures where in the
\emph{observed event stream} a seat's declares landed (mean declare event index over the
final observed index), which is the observable stand-in for the decline-counting
``declare move index'' a god's-eye harness would use.

This blind spot matters more than it looks, and Section~\ref{sec:limit} returns to it: some
of the very decisions on which styles differ are declines, and those decisions leave no trace
at all.

\section{The feature set}
\label{sec:features}

\begin{table}[t]
\centering
\caption{The 14 features of the shipped observation vector, in classifier order
\cite{repo}. All are rates or shares; denominators are floored at 1.}
\label{tab:features}
\small
\begin{tabular}{@{}llp{8.0cm}@{}}
\toprule
\# & Feature & Definition \\
\midrule
1 & \code{hitRate} & hits / asks \\
2 & \code{askDiversity} & distinct books asked into / asks \\
3 & \code{sameBookRepeatRate} & consecutive same-book asks by the seat / (asks $-$ 1) \\
4 & \code{certainAskShare} & asks whose named card was publicly located at the target / asks
    (a guaranteed public hit) \\
5 & \code{deadAskShare} & asks whose named card was publicly located elsewhere / asks
    (a guaranteed public miss) \\
6 & \code{leakyAskShare} & asks into books where the asker's team already publicly accounted
    for $\geq 4$ of six cards (located plus certified, never double-counted) / asks \\
7 & \code{completionAskShare} & asks into books with five cards publicly located at the
    asker's team / asks \\
8 & \code{missFewestShare} & of the seat's misses, the fraction whose target held strictly
    the fewest cards among live opponents (replayed counts at ask time) \\
9 & \code{missMostShare} & \ldots{} strictly the most \\
10 & \code{askShare} & the seat's asks / all observed events \\
11 & \code{declareShare} & the seat's declares / all observed events \\
12 & \code{foreignDeclareShare} & foreign declares / declares (exact; Section~\ref{sec:exact}) \\
13 & \code{ownHandOnlyShare} & own-hand-only declares / declares (exact) \\
14 & \code{declareBackload} & mean declare event index / final observed event index \\
\bottomrule
\end{tabular}
\end{table}

Table~\ref{tab:features} lists the vector. Three design decisions deserve their measured
justifications.

\paragraph{Shares, never raw counts.} Every entry is a rate or a share, length-invariant by
construction. This is a correction, not a preference: the first calibration ran over raw
counts and read every short cross-play game as ``not the style'' simply because the mirror
games it was calibrated on ran longer --- a fact about the \emph{opponents}, not the seat
\cite{repo}. A count observed at 150 events is incommensurable with a fingerprint calibrated
at a different game length; a share projects a 70-event clinch and a 200-event grind into the
same space. Checkpoint buckets (Section~\ref{sec:classifier}) still matter on top, because
shares themselves drift over a game: declares concentrate late, and public certainty
accumulates.

\paragraph{Deliberate exclusions.} Raw \code{asks}, \code{hits} and \code{declaresCorrect}
are absent. The first two are already the denominators inside \code{hitRate} and
\code{askShare}; correctness is near ceiling for every full-strength style (roster declare
precision $0.946$--$0.990$ \cite{v05}), so it separates nothing and would only add
correlated noise to a diagonal model.

\paragraph{A feature that never fired.} In the entire calibration --- 8{,}100 seat-vectors
from 1{,}350 mirror games --- \code{completionAskShare} is identically zero for all nine
styles \cite{fingerprints}. The interpretation (offered as interpretation) is structural: a
full-strength engine whose team publicly holds five cards of a book almost always has the
certainty to declare it rather than ask into it. Within this roster the column separates
nothing and, under the classifier's variance floor, penalises all styles equally; it is
retained because it is aimed at play outside this roster --- weaker or human play, where
chasing a publicly finished book is a real behaviour. Relatedly, \code{deadAskShare} is the
public signature of the contained-book turn-pass studied in the second companion
\cite{containedpaper}: that mechanism's deliberate guaranteed misses are exactly what
feature~5 counts.

\section{Calibration and the classifier}
\label{sec:classifier}

\subsection{Fingerprints from mirror games}

Fingerprints are generated offline and committed \cite{fingerprints}: per style, 150 \usff{}
mirror games (all six seats the same style), seeds \code{fingerprint-v1-<style>-<i>}, start
seat $i \bmod 6$, observed through the \emph{same code path} the classifier uses --- the
generator and the classifier share the extractor, or the fingerprints would describe a
different instrument than the one that reads them \cite{repo}. Each game contributes six
seat-vectors, so the full-log bucket holds 900 vectors per style.

Because several behaviours drift over a game, each style carries one fingerprint per
\emph{checkpoint bucket}: log prefixes of 60, 120, 200 and 300 events, plus \code{full}
(completed games). A bucket only collects games whose log outlived it, and the committed
provenance records how quickly the population thins: at 120 events, between 180 (Blitz) and
882 (Turtle) of the 900 seat-vectors remain; only the Turtle's games reach 200 events in any
number (246 vectors); \emph{no} game reached 300, so that bucket carries the full-game
statistics with a recorded sample count of zero \cite{fingerprints}. Bucket selection at read
time: a view whose log ends in \code{game\_over} reads \code{full} --- it \emph{is} a
completed game, whatever its length --- otherwise the nearest checkpoint at or below the
observed event count, falling back to 60 below 60 events \cite{repo}.

\subsection{The model, and two corrections made by measurement}
\label{sec:corrections}

The model is deliberately the simplest thing that can be honest: a diagonal Gaussian per
style over the 14-vector. For an observed vector $x$ and style $s$ with per-feature
calibration mean $\mu_{s,i}$ and standard deviation $\sigma_{s,i}$,
\[
\ell_s(x) \;=\; -\tfrac12 \sum_{i=1}^{14} z_{s,i}^2,
\qquad
z_{s,i} \;=\; \frac{x_i - \mu_{s,i}}{\tilde\sigma_{s,i}},
\]
and a softmax over the nine $\ell_s$ gives the raw posterior. Diagonal on purpose: with 14
features and a few hundred vectors per bucket, a full covariance would be an estimate of
noise, and accuracy is \emph{measured} downstream (Section~\ref{sec:accuracy}), never assumed
here. Two modelling choices were forced by calibration failures \cite{repo}:

\begin{enumerate}
\item \textbf{The Gaussian normaliser is dropped.} The textbook log-likelihood carries a
  $-\ln\sigma_{s,i}$ term per feature. With sample SDs over a few hundred vectors,
  $\sum_i -\ln\sigma_{s,i}$ is a large \emph{constant} bonus for whichever style happened to
  calibrate tightest --- and in the first calibration it handed the Ghost 40 of 60
  mirror-Balanced reads. Without the term the score is a pure $z$-distance, whose argmax at a
  style's own calibration mean is that style itself. The cost is that the model is no longer
  a normalised density; the benefit, measured, is that it stops rewarding calibration
  tightness as if it were evidence.
\item \textbf{Per-style SDs are shrunk halfway to pooled:}
  $\tilde\sigma_{s,i}^2 = (\sigma_{s,i}^2 + \bar\sigma_i^2)/2$, where $\bar\sigma_i$ is the
  root-mean-square of the nine styles' SDs for feature $i$. Raw per-style SDs make the
  widest-variance style an outlier magnet --- everything unusual ``must be the Scout'' ---
  while fully pooled SDs discard real information (the Turtle genuinely is more repetitive
  than its mean alone says). The halfway blend beat both endpoints on cross-play calibration
  probes and tied them on mirror self-reads.
\end{enumerate}

Both corrections are small, empirical, and the kind of thing a methods section usually omits.
They are reported here because each one, left unfixed, produces a confidently wrong
instrument, and because neither is specific to this game: any few-sample per-class Gaussian
over behavioural rates faces both.

\subsection{Honesty mechanisms}

Two further mechanisms are structural rather than statistical \cite{repo}. \emph{Sample-size
damping}: the posterior is blended toward uniform by $w = \min(1, \mathrm{asks}/12)$, so at
zero observed asks the answer is exactly ``I don't know,'' and full confidence is reachable
only once a seat has asked a dozen times --- asks, not events, because asks are the seat's
own observed choices, and a seat that has barely moved in a long game is still unread.
\emph{A variance floor} of $0.01$ prevents the one literally-constant calibration column
(Section~\ref{sec:features}) from turning a single stray count into an infinite log-penalty;
it becomes a very strong finite vote instead.

The unit-test suite pins the contract rather than the accuracy: the posterior is a proper
distribution; zero asks damp to exactly uniform; the bucket rule is pinned point by point;
the classifier is deterministic and never mutates a frozen view; and one deliberately loose
real-game smoke test asserts only better-than-chance on the roster's most separable axis
(Turtle-family reads at roughly $0.31$ against a $2/9 \approx 0.22$ chance rate)
\cite{repo}. Real accuracy is a measured number (Section~\ref{sec:accuracy}), not a unit-test
assertion --- a tight bound there would be a number the suite cannot honestly own.

\section{Separability: the roster's geometry in observable space}
\label{sec:separability}

\begin{table}[t]
\centering
\caption{Pairwise distances between the committed full-log fingerprints
\cite{fingerprints}, in pooled-$z$ units: $d(a,b) = (\sum_i ((\mu_{a,i} -
\mu_{b,i})/\bar\sigma_i)^2)^{1/2}$ with $\bar\sigma_i$ the RMS of the nine styles'
calibration SDs for feature $i$ (floored at $0.01$; the floor binds only on the identically
zero completion column). Computed directly from the committed table; reproducible with a
few lines of \code{node}.}
\label{tab:dist}
\footnotesize
\begin{tabular}{@{}lcccccccc@{}}
\toprule
 & Balanced & Blitz & Punter & Banker & Turtle & Hoarder & Scout & Ghost \\
\midrule
Blitz     & 0.22 &      &      &      &      &      &      &      \\
Punter    & 0.24 & 0.25 &      &      &      &      &      &      \\
Banker    & 0.20 & 0.23 & 0.24 &      &      &      &      &      \\
Turtle    & 3.22 & 3.35 & 3.18 & 3.28 &      &      &      &      \\
Hoarder   & 1.74 & 1.79 & 1.74 & 1.76 & 2.81 &      &      &      \\
Scout     & 1.66 & 1.79 & 1.64 & 1.69 & 2.12 & 1.97 &      &      \\
Ghost     & 1.02 & 1.04 & 1.15 & 0.97 & 3.70 & 2.15 & 2.07 &      \\
Archivist & 0.95 & 1.08 & 0.96 & 0.99 & 2.58 & 1.74 & 0.75 & 1.42 \\
\bottomrule
\end{tabular}
\end{table}

Table~\ref{tab:dist} is the paper's central picture: the distance between every pair of
committed full-log fingerprints, in pooled-$z$ units. Four structures stand out, and each has
a mechanism.

\paragraph{The quadrangle.} Balanced, Blitz, Punter and Banker sit pairwise
$0.20$--$0.25$ apart --- over a fourteen-dimensional space, nearly a single point. This is
the inert-axis finding of the companions \cite{inertpaper,v05} arriving in observable space.
Those four styles are exactly the ones whose defining knob was the declare threshold, an axis
measured to select identical declares across the roster's entire range; what behavioural
difference survives (their ask-scoring weights) amounts to $0.39$--$1.19\%$ of decisions
against the control \cite{styles}, and $\le 1.3\%$ of decisions leaves almost no trace in
fourteen public rates. The tournament separates these styles by up to $5.4$ points of mean
score \cite{v05}; the public log, in one game, essentially cannot.

\paragraph{The far pole.} The Turtle sits $3.2$--$3.7$ from the quadrangle and the Ghost, and
no closer than $2.1$ to anything --- the roster's one style an observer can genuinely expect
to spot. Its per-feature signature is broad, not narrow: relative to Balanced,
\code{missFewestShare} $+1.7\bar\sigma$, \code{leakyAskShare} $+1.6\bar\sigma$,
\code{deadAskShare} $+1.4\bar\sigma$, \code{hitRate} $-0.9\bar\sigma$, \code{askDiversity}
$-0.8\bar\sigma$. The mechanism is its own-hand-only declare gate \cite{v05}: refusing to
bank split books keeps resolved information on the table, so its asks are measured against a
position the whole table has already solved.

\paragraph{The Hoarder, and what the hoard-gate fix bought.} The Hoarder sits $1.74$ from
Balanced --- and essentially the whole distance is one axis, \code{foreignDeclareShare} at
$+1.5\bar\sigma$ (calibration mean $0.225$ against the control's $0.009$). This is the
audited hoard gate's measured side effect: a style that refuses declares which spend its
licences ends up declaring disproportionately the books that cost it nothing, i.e.\ foreign
ones \cite{styles,v05}. Two facts qualify the $1.74$. First, it exists only because of the
audit: the first Hoarder implementation was byte-identical to Balanced --- zero divergent
decisions in 275{,}380 --- so its observable distance was exactly $0$ \cite{v05}. A style
must first \emph{exist} to be observed. Second, the distance is strongly time-dependent: at
the 60-event bucket the Hoarder--Balanced distance is $0.15$, because declares --- the only
place its signature lives --- concentrate late. Early-game reads of the Hoarder are close to
structurally impossible.

\paragraph{Scout $\approx$ Archivist.} The two information styles sit $0.75$ apart, with no
single feature differing by more than $0.44\bar\sigma$. The Archivist is the interesting
half of the pair: its name-giving move --- declaring books it holds no card of, for its
teammates --- fires at a shipped threshold of $0.95$ and almost never triggers, so its
measured foreign-declare rate ($0.007$) is actually \emph{below} the control's ($0.014$)
\cite{v05}. Its public signature is therefore not its thesis but its ask-weight side
effects, which resemble the Scout's. A style whose defining mechanism rarely fires is, to
this channel, mostly its side effects.

\section{Measured accuracy}
\label{sec:accuracy}

\subsection{Design}

Accuracy is measured on games the calibration never saw, in a setting deliberately harder
than calibration: \emph{cross-play}. All 36 unordered style pairings play 50 single games
each (team~0 one style, team~1 the other; seeds \code{clsacc-v1-<i>}, disjoint from the
calibration prefix; start seat rotating), for 1{,}800 games. Each game is classified from log
prefixes of 40, 80, 150 and 250 events and from the full log, using the committed classifier
and fingerprints end to end; a truncation is scored only for games whose log outlived it.
Each checkpoint therefore yields up to $1{,}800 \times 6 = 10{,}800$ seat-reads, each scored
top-1 against the seat's true style among nine candidates (chance $1/9 \approx 11.1\%$)
\cite{artifact}. The distribution shift from mirror calibration to cross-play deployment is
intentional --- it is the shift the adaptive engine \cite{v10} actually faces --- and its
cost is visible in the numbers below.

\subsection{The accuracy curve}

\begin{table}[t]
\centering
\caption{Single-game top-1 accuracy over nine styles, by observed-event checkpoint
\cite{artifact}. Chance is $1/9 \approx 0.111$. A checkpoint scores only games whose log
outlived it; the 150-event row's small, style-skewed survivor population (112 of 1{,}800
games, dominated by the long-game styles) is why its higher figure is not comparable to the
others, and no game reached 250 events.}
\label{tab:curve}
\small
\begin{tabular}{@{}lccc@{}}
\toprule
Checkpoint (events) & Seat-reads & Top-1 & vs.\ chance \\
\midrule
40         & 10{,}800 & 0.164 & $1.5\times$ \\
80         & 10{,}464 & 0.202 & $1.8\times$ \\
150 (survivors only) & 672 & 0.375 & --- \\
full log   & 10{,}800 & 0.224 & $2.0\times$ \\
\bottomrule
\end{tabular}
\end{table}

Table~\ref{tab:curve}: accuracy rises from $16.4\%$ at 40 events through $20.2\%$ at 80 to
$22.4\%$ on the full log --- twice chance, and no better than that. The curve's shape is the
damping and the buckets doing their jobs (early reads are pulled toward uniform; later reads
have more declares to work with), and its \emph{level} is the separability geometry of
Section~\ref{sec:separability} made concrete. One row needs its caveat stated plainly: the
$37.5\%$ at 150 events is computed over the 672 seat-reads of the 112 games that lasted that
long, a population heavily skewed toward the long-game styles (Turtle 186 and Scout 180 of
the 672 seats), which are also the separable ones. It is a fact about survivors, not a point
on the same curve.

\subsection{Per-style accuracy and the confusion structure}

\begin{table}[t]
\centering
\caption{Full-log confusion, row-normalised to percent (1{,}200 seat-reads per row)
\cite{artifact}. Rows are true styles, columns predictions; the diagonal is bold.}
\label{tab:confusion}
\footnotesize
\setlength{\tabcolsep}{4.5pt}
\begin{tabular}{@{}lccccccccc@{}}
\toprule
True $\backslash$ read & Bal & Bli & Pun & Ban & Tur & Hoa & Sco & Gho & Arc \\
\midrule
Balanced  & \textbf{5.6}  & 8.3 & 15.2 & 14.4 & 2.2  & 4.4  & 15.3 & 27.0 & 7.8 \\
Blitz     & 5.3 & \textbf{9.3}  & 16.6 & 14.2 & 1.3  & 4.8  & 14.3 & 26.1 & 8.2 \\
Punter    & 4.8 & 9.1 & \textbf{12.9} & 12.5 & 3.2  & 5.2  & 16.4 & 27.5 & 8.4 \\
Banker    & 5.5 & 8.5 & 14.5 & \textbf{12.6} & 2.8  & 3.4  & 16.1 & 32.3 & 4.3 \\
Turtle    & 2.2 & 2.8 & 4.4  & 7.7  & \textbf{27.2} & 7.8  & 23.1 & 19.3 & 5.7 \\
Hoarder   & 1.2 & 3.7 & 6.3  & 9.0  & 9.3  & \textbf{30.9} & 17.3 & 18.3 & 3.9 \\
Scout     & 2.5 & 4.1 & 4.1  & 4.9  & 15.2 & 1.2  & \textbf{41.3} & 16.9 & 9.9 \\
Ghost     & 3.2 & 5.8 & 11.9 & 11.4 & 0.8  & 2.2  & 10.3 & \textbf{50.4} & 3.9 \\
Archivist & 2.8 & 4.8 & 6.0  & 5.6  & 12.3 & 1.9  & 33.1 & 21.8 & \textbf{11.8} \\
\bottomrule
\end{tabular}
\end{table}

The headline $22.4\%$ decomposes into two regimes (Table~\ref{tab:confusion}). The styles
with a real observable signature are read far above chance: Ghost $50.4\%$, Scout $41.3\%$,
Hoarder $30.9\%$, Turtle $27.2\%$. The quadrangle is read at $5.6$--$12.9\%$ ---
\emph{at or below} chance for Balanced and Blitz --- exactly as the distance table predicts
for four styles $0.2$ apart under calibration SDs an order of magnitude larger.

This split was pre-stated. The run's committed predictions include (\code{P4}) that
accuracy ``should be good for turtle, ghost and hoarder,'' and heavily confused inside the
Balanced--Blitz--Punter--Banker quadrangle, on the Stage-1B separability read
\cite{artifact}. The recorded verdict is \emph{mixed}, and honestly so: the quadrangle
confusion held (mean $10.1\%$ against a $\le 35\%$ condition), but the separable trio
under-shot its $\ge 50\%$ bar (mean $36.2\%$; only the Ghost cleared it) --- and quadrangle
errors did \emph{not} mostly stay inside the quadrangle ($30$--$40\%$ against a $\ge 50\%$
condition). Where they went instead is the table's most interesting structure.

\paragraph{Ghost and Scout as attractors.} Summing columns: the Ghost is predicted for
$2{,}875$ of $10{,}800$ reads ($26.6\%$) and the Scout for $2{,}245$ ($20.8\%$) ---
together $47.4\%$ of all reads, for two styles that are $2/9$ of the truth. Every quadrangle
style's single most common read is Ghost ($26$--$32\%$); the Archivist is read as Scout
nearly three times as often as as itself ($33.1\%$ vs.\ $11.8\%$).

We offer an interpretation, and flag it as interpretation rather than measurement. The
per-feature decompositions of Section~\ref{sec:separability} show the roster's two strongest
observable axes are the ones the Ghost and the Scout respectively own: the
ask-repetition/hit-rate direction (Ghost: fewer same-book repeats, more certain asks, higher
hit rate --- leak avoidance's public shadow) and the dead-ask/targeting direction (Scout:
public misses and miss-target structure --- information buying's public shadow). A noisy
seat-vector that wanders from the quadrangle's shared centre must wander \emph{in some
direction}, and these two poles are the nearest distinct means in the two highest-variance
observable directions; under a pure $z$-distance argmax, noise falls toward them. On this
reading the attractor structure is a property of the observable space's geometry, not a
defect of the model --- but distinguishing those two explanations would need a calibration
ablation this paper does not contain.

\section{The fundamental limit, and the program this instrument belongs to}
\label{sec:limit}

\subsection{The divergence ceiling}

The decision-divergence instrument of the v0.5 paper \cite{v05,styles} measures, by
replaying identical views and seeds, the exact fraction of decisions on which each style
differs from the Balanced control: Scout $2.89\%$, Ghost $2.15\%$, Archivist $1.66\%$,
Banker $1.19\%$, Turtle $0.83\%$, Punter $0.47\%$, Hoarder $0.47\%$, Blitz $0.39\%$
\cite{styles}. Those same measurements put the decision density at roughly $675$ policy
decisions per game across all six seats (declare-window decisions included), i.e.\ on the
order of $110$ per seat.

The arithmetic of the ceiling follows directly. A Blitz seat makes about $0.4$
Blitz-revealing decisions per game; a Scout seat about $3$. Every other decision is, by
construction, \emph{identical to the control's} --- carrying zero information about the
style. No observer of any channel, public or god's-eye, can classify a seat from decisions
the seat's style never influenced; one game of a $0.39\%$-divergent style simply does not
contain the evidence, and the measured quadrangle accuracy of $5.6$--$12.9\%$ should be read
as that fact, not as a classifier deficiency. The public channel then tightens the bound
further: a share of the divergent decisions are declare-window choices, and a divergent
\emph{decline} emits no event at all (Section~\ref{sec:blind}) --- some of the few
style-revealing decisions are invisible in principle. Styles are distinguishable in one game
roughly in proportion to where their divergence lands: the Turtle's $0.83\%$ is highly
visible because it lands on asks and declares that reshape the whole public position, while
the quadrangle's $0.39$--$1.19\%$ lands on ask-scoring margins that the log renders as the
same kind of event either way.

The constructive reading of the ceiling is about aggregation: at one game the observer is
evidence-starved, not model-starved, so the natural upgrades are longitudinal ---
accumulating observations across games of a persistent opponent --- rather than
architectural. That is the direction the synthesis's program points, and it is future work.

\subsection{The detection-panel program}

The project's research synthesis \cite{playstyles} specifies (its \S VI.4) an
opponent-detection statistics panel --- every statistic computable from the public log alone
--- as the feature set of its cross-game opponent-profiling style, ``The Reader'' (its S35).
This paper's instrument is best understood as implementing and \emph{measuring} a slice of
that panel against a concrete roster. Shipped, with the panel's intended detection target:
the dead-ask rate (its known-negative-ask detector), the certain-ask rate (its
confirmation-ask / claim-imminent detector), ask diversity with hit rate (its signal-broker
row), the same-book repeat rate (its chain-discipline row), the miss-target shares (the
observable core of its turn-parking and targeting rows), declare timing via
\code{declareBackload} (its claim-delay row, on the observable stand-in of
Section~\ref{sec:blind}), and the exact foreign / own-hand-only declare shares (its claimer
taxonomy).

Unshipped, and honestly labelled future work: the blackballing detectors (collapse of a
seat's inbound-target share); the evidence-age decay of hit rate --- the synthesis's
candidate best single skill estimator --- which needs per-ask evidence dating; the possum
detector (deducible certainties high while certain asks stay near zero), which requires
running the full deduction engine over the log and is deliberately outside this observation
layer's no-inference scope; claim-threshold estimation from provably contained books; and
the entire convention-detection family --- the Salahuddin-convention detector and the
permutation-null excess-mutual-information test --- which has nothing to detect here, since
this roster plays no conventions, and which additionally needs revealed-hand data retained
across games. None of the unshipped rows would lift the Section~\ref{sec:limit} ceiling for
this roster in a single game; the panel's value concentrates in the cross-game, human-facing
setting the synthesis designed it for.

\section{Honest scope}
\label{sec:scope}

Everything above is conditional, and the conditions are worth stating as bluntly as the
results.

\begin{enumerate}
\item \textbf{This roster, this rule set.} Features were designed and fingerprints calibrated
  for the nine committed styles under \usff{} as pinned by the rules hash
  \cite{artifact,rules}. The feature set was written by authors who knew the roster ---
  \code{foreignDeclareShare} exists because foreign-declaring styles do --- so accuracy on a
  novel, adversarial, or human style is unmeasured, and the design leakage means it cannot be
  extrapolated from Table~\ref{tab:confusion}.
\item \textbf{Full-strength shared inference.} Every observed seat runs the identical
  full-strength inference engine \cite{v05}. Behavioural variation from \emph{skill} ---
  memory limits, deduction errors --- is entirely absent from calibration. Human play varies
  along exactly that axis; it is out of distribution, and the synthesis's panel rows for
  skill estimation (evidence-age decay) are the missing instrument, not this one.
\item \textbf{Mirror calibration, cross-play deployment.} Fingerprints come from mirror
  games; accuracy is measured on cross-play, where game length and opponent behaviour differ.
  This shift is intentional (it is the deployment condition) but it means the reported
  accuracy is neither a ceiling nor a floor for other populations.
\item \textbf{Single-game observation only.} Nothing here aggregates across games. The
  divergence arithmetic of Section~\ref{sec:limit} says single-game reads of near-control
  styles are evidence-starved in principle; cross-game profiles are the synthesis's S35
  program and remain future work.
\item \textbf{Deterministic policies.} All variation comes from seeds; a stochastic or
  deliberately deceptive policy (the synthesis's possum) faces detectors this layer
  deliberately does not include.
\end{enumerate}

\section{Conclusion}
\label{sec:conclusion}

The \usff{} public log is an unusually generous observation channel --- it locates cards,
certifies licences, and reveals declaration truth --- and this paper measured how much play
style survives the trip through it. The answer has a clean shape. What the channel certifies,
it certifies exactly: foreign and own-hand-only declares are noiseless observations, and a
14-feature share vector over the log supports a simple, honest classifier once two measured
corrections (no normaliser; half-pooled variances) are applied. What the channel recovers is
governed by geometry that the roster's own companion findings predict: the four styles whose
labelled axis is inert collapse to a point $0.20$--$0.25$ wide and read at or near chance,
while the styles whose divergence lands on public, position-reshaping events --- Turtle,
Hoarder, Scout, Ghost --- are read at $27$--$50\%$ against $11.1\%$ chance. And what limits
the whole enterprise is not the model but the material: styles that diverge on
$0.39$--$2.89\%$ of decisions offer an observer one to three revealing decisions per seat per
game, some of them structurally invisible declines. Twice chance at one game is what this
channel, this roster, and honest calibration yield; the road past it runs through cross-game
aggregation, not a better single-game model. Whether even a perfect read would be worth
anything to an adapting agent is the companion paper's question --- and its answer belongs to
that paper, not this one \cite{v10}.

\medskip
\noindent\textbf{Reproducibility.} The engine and classifier are deterministic. The committed
fingerprint table carries its full provenance (generator command, seed prefix, per-bucket
sample counts) \cite{fingerprints}; the accuracy experiment's committed artifact carries the
run configuration, health counters (all clean), records digest and engine commit
\cite{artifact}; and the distance matrix of Table~\ref{tab:dist} recomputes from the
committed fingerprints alone.

\section{Addendum: re-measured on the regenerated artifact}
\label{sec:addregen}

\emph{September 2026.} Two things changed under this paper after it was written, and the
classifier was re-calibrated and re-measured on both. The first is the concession layer of
v2.0: every roster style now carries \code{defuse} $=1$, which changes what the nine styles
\emph{do} and therefore what a public log about them contains. The second is a correction to
\code{RULES\_US54.md} \S4, where a rule marked \code{[DERIVED]} advanced the turn to the next
seat holding cards when a declare emptied the turn-holder --- an opponent, where the owner's
rulebook passes to a teammate \cite{rules}. The fingerprints of Section~\ref{sec:classifier}
were regenerated from fresh mirror games and the accuracy run of Section~\ref{sec:accuracy}
re-played on the same 1{,}800 held-out cross-play games \cite{artifact}.

\textbf{The headline is unchanged and the shape of the curve with it.}

\begin{table}[t]
\centering
\caption{Single-game top-1 over nine styles, as published and as regenerated
\cite{artifact}. Chance is $11.1\%$ throughout.}
\label{tab:addcurve}
\footnotesize
\begin{tabular}{@{}lrr@{}}
\toprule
log horizon & as published & regenerated \\
\midrule
40 events   & $16.4\%$ & $16.1\%$ \\
80 events   & $20.2\%$ & $21.7\%$ \\
full log    & $22.4\%$ & $22.0\%$ \\
\bottomrule
\end{tabular}
\end{table}

Twice chance and no better, at every horizon, on both runs. Nothing in
Sections~\ref{sec:limit} or \ref{sec:scope} moves: the ceiling is still the material rather than
the model, and it is still set by how few decisions a style reveals.

\subsection{What did move: which style is hardest to read}

The per-style column is a different matter, and this addendum will not smooth it over.
As published, end-of-game accuracy ran from $50.4\%$ on the Ghost down to $5.6\%$ on the Balanced
control, and the Balanced-is-invisible result was used in Section~\ref{sec:limit} as the sharpest
illustration of the divergence argument. \textbf{On the regenerated artifact the floor belongs to
a different style.}

\begin{table}[t]
\centering
\caption{End-of-game top-1 by style on the regenerated artifact, 1{,}200 seat-reads each
\cite{artifact}. The published run's two anchors were Ghost $50.4\%$ and Balanced $5.6\%$.}
\label{tab:addbystyle}
\footnotesize
\begin{tabular}{@{}lrlr@{}}
\toprule
style & top-1 & style & top-1 \\
\midrule
Ghost     & $41.4\%$ & Blitz     & $13.8\%$ \\
Scout     & $32.7\%$ & Banker    & $11.3\%$ \\
Turtle    & $32.1\%$ & Archivist & $10.6\%$ \\
Hoarder   & $29.3\%$ & Punter    & $\mathbf{3.6\%}$ \\
Balanced  & $23.3\%$ &           &           \\
\bottomrule
\end{tabular}
\end{table}

The Ghost is still the most legible style and still by a wide margin. But Balanced has risen from
$5.6\%$ to $23.3\%$ and \textbf{Punter has fallen to $3.6\%$}, so the hardest style to identify
from a public log is now the roster's strongest one rather than its control.

Two readings are available and this paper commits to neither. The mechanical one is that the
control is no longer the roster's midpoint: a mechanism every style inherits does not move a
style relative to the others, but it does move where the \emph{classifier's} decision boundaries
fall, and Section~\ref{sec:corrections}'s two calibration corrections were both fitted against a
population that no longer exists. The substantive one is that Punter's behaviour has become the
population's centre of mass, which would make it the style an unsure classifier defaults
\emph{away} from. \textbf{Distinguishing them needs a divergence measurement of the regenerated
roster against its control, and that has not been run.}

\subsection{What this addendum does not re-derive}

The separability geometry of Section~\ref{sec:separability} --- the cluster widths, the Turtle's
distance from the inert group --- and the two calibration corrections of
Section~\ref{sec:corrections} were not recomputed. They are statements about the feature space of
the pre-concession roster, they are reported here as such, and the $40$ of $60$ control-game reads
that forced the normaliser correction is a fact about that calibration rather than about this one.
Anyone re-deriving them should expect the numbers to move by at least as much as
Table~\ref{tab:addbystyle} did.

\begin{thebibliography}{9}
\small

\bibitem{v05}
A.~Hsieh.
\newblock FishAI v0.5: Measuring play style without skill in a 54-card Literature variant.
\newblock FishAI project; companion paper, published alongside this one, 2026.

\bibitem{v10}
A.~Hsieh.
\newblock FishAI v1.0: When best-response adaptation is worth less than nothing.
\newblock FishAI project; companion paper, 2026.

\bibitem{inertpaper}
A.~Hsieh.
\newblock The inert axis: when a style parameter is wired, swept, and never reached.
\newblock FishAI project; companion paper, 2026.

\bibitem{containedpaper}
A.~Hsieh.
\newblock The contained book: an absorbing resource measured to be worth nothing.
\newblock FishAI project; companion paper, 2026.

\bibitem{playstyles}
FishAI project.
\newblock PLAYSTYLES.md --- Literature / Canadian Fish: play styles, rule dialects, and
engine specification.
\newblock Internal research synthesis with per-claim evidence tiers, 2026. \S VI.4 specifies
the opponent-detection statistics panel; S35 (``The Reader'') the cross-game profiling
program; S21 and S36 the possum and Salahuddin-convention detection targets.

\bibitem{rules}
FishAI project.
\newblock RULES\_US54.md --- the ``US student'' 54-card variant.
\newblock Repository rules specification with derived consequences and test vectors, 2026.
Rows 6--7 (ask licences), row 9 (hit transfer), and the declare-window mechanics cited in
Section~\ref{sec:channel}.

\bibitem{styles}
FishAI project.
\newblock STYLES.md --- the play-style roster under \usff{}.
\newblock Repository document; \S3.1.1 the Hoarder inertness audit, \S6.1 the inert
declare-threshold axis and per-style decision-divergence rates, 2026.

\bibitem{fingerprints}
FishAI project.
\newblock Committed fingerprint calibration \code{lib/engine/bots/data/fingerprints.ts}.
\newblock Generated by \code{node scripts/gen-fingerprints.mjs --games 150}: 150 \usff{}
mirror games per style, seed prefix \code{fingerprint-v1}, checkpoint buckets
$\{60, 120, 200, 300, \code{full}\}$ with per-bucket sample provenance; generated
2026-08-24.

\bibitem{artifact}
FishAI project.
\newblock Adaptive-lab artifact \code{src/lab/data/adaptive-results.json}.
\newblock Classifier-accuracy block: 36 cross-play style pairings $\times$ 50 games (seed
prefix \code{clsacc-v1}), truncation checkpoints $\{40, 80, 150, 250\}$ plus full log;
records digest \code{ce3316641a8f38fe}; engine commit \code{45ec9f3} (recorded
\code{45ec9f3-dirty}); health counters all zero; generated 2026-08-24.

\bibitem{repo}
FishAI project.
\newblock FishAI repository: engine, bots, simulation laboratory, and results site.
\newblock \code{lib/engine/bots/observe.ts} (the observation layer and its header
derivations), \code{lib/engine/bots/classify.ts} (the classifier and the two calibration
corrections), \code{tests/bots/observe.test.ts} and \code{tests/bots/classify.test.ts} (the
pinned contracts), 2026.

\end{thebibliography}

\end{document}
